\section{Discussion and conclusion}\label{sec:discussion}

The methodology of Section~\ref{sec:method}, together with the crossover map of
Section~\ref{sec:crossover}, is intended to make one specific question answerable in
a neutral, reproducible way: \emph{what resilience does a given sensor architecture
buy, and over what operating envelope?} Rather than asserting a verdict, Kshana
\citep{baweja2026kshana} fixes a scenario, a seed, and an engine version, routes the
quantum and classical candidates through a single code path, and labels every
figure-of-merit as externally validated or merely modelled. The output of that
discipline is not a marketing number but an envelope: a region of the operating space
in which one architecture dominates, a region in which it does not, and a break-even
boundary between them carried with a $95\%$ confidence band.

This reframing matters because the comparison it supports is different from the one the
literature usually runs. The question \emph{``is quantum sensing better?''} is already
well populated by device papers and surveys
\citep{templier2022,coldatomnavreview2022,quantumpntreview2026}; it is answered
affirmatively on isolated axes and, on its own, settles little for a system designer.
The crossover map asks instead \emph{``where does the quantum lever change the
outcome?''}, and answers it openly, with the uncertainty made explicit rather than
hidden. The inertial study locates that boundary as a vibration tolerance that grows
with outage duration --- the break-even platform-vibration PSD rises from the order of
$10^{-5}$ at a $30$~s outage to the order of $10^{-4}$~(m/s$^2$)$^2$/Hz near a
$1800$~s outage --- so the cold-atom advantage is real for long, quiet holdovers and
erodes under aggressive platform vibration, exactly where a TRL-honest answer should
place it. The clock study is starker still: the deployed chip-scale reference crosses a
$1\,\mu$s timing budget at roughly $3.0\times10^{4}$~s ($\sim$8.5~h) of holdover, while
the ground-lab optical-lattice and flight-qualified maser curves stay below that budget
across the full $24$~h sweep. These are performance-model projections from cited device
coefficients, not measurements, and they are TRL-tagged as such; the contribution is the
reproducible map and its break-even contour, not any single point on it.

Interoperability is what turns this from a closed demonstrator into an on-ramp. Because
the engine reads and writes the standard exchange formats (RINEX, SP3, CCSDS OEM/OMM),
a crossover study or a propagated state can move into established pipelines --- RTKLIB
\citep{takasu2009}, GMAT \citep{nasa_gmat}, Orekit \citep{orekit} --- rather than
living only inside Kshana. The single integrated stack runs natively, from Python, and
in the browser via WebAssembly; in-browser SGP4 is not itself new
\citep{satellitejs,sgp4crate}, and the delta is only that one engine runs the full
integrated stack with zero install, lowering the cost of reproducing a result to opening
a page.

Several limitations defined in Section~\ref{sec:limitations} double as the natural work
ahead. The PVT path is scenario-driven rather than fed by real observations; ingesting
real GNSS observations to close a measurement-level loop is the most consequential next
step. The spoofing detector is characterised against published TEXBAT
\citep{humphreys2012texbat} parameters, not raw IQ; ingesting raw IQ and reporting a
detection ROC would move it from the not-externally-validated tier into a measured one.
The force-model validation by ephemeris fitting reported in
Section~\ref{sec:validation} is single-arc and single-satellite; extending it to
multi-satellite, multi-day fits would tighten the residual envelope and test the static
drag, cannonball-SRP, and Schwarzschild defaults under more diverse dynamics. Finally,
the crossover map presently sweeps published-parameter envelopes one lever at a time;
propagating the full coefficient uncertainty jointly --- and adding same-TRL pairings
across all three levers --- would turn today's three studies into a single crossover
atlas of resilient PNT.

What we offer, then, is less a claim of superiority than a way of arguing about it
honestly. The agency-orbit work is a force-model validation by ephemeris fitting, not
orbit determination; the headline full-arc Galileo residual is $0.61$~m at degree and
order $70$, with the $0.13$~m and $0.10$~m reduced-dynamic figures stated as
empirical-acceleration-tier bounds rather than measures of force-model accuracy; the
LRO residual is $6.6$~m, above the $5$~m target, reported as such and reproduced as a
known result \citep{bertone2021grail,li2023leopod} against a DE-grade ANISE oracle
\citep{nyxanise} that we credit rather than claim to beat; and every quantum figure is a
projection, every spoofing claim a parameter-level characterisation. To our knowledge no
open tool runs this full integrated PNT-resilience stack with these honesty controls
attached. That stance --- validated-or-not-modelled stated per figure-of-merit, the
quantum-versus-classical axis swapped through one neutral code path, and every headline
number traceable to a committed scenario, seed, and test --- is the differentiator we
mean to defend, consistent with the reproducibility tradition this work builds on
\citep{fernandezprades2018reproducibility,nasastd7009}.
